\paragraph{A uniform periodic extension of the arithmetic divisor.}
The following statement extends the exact finite-field readout argument.
It is not a theorem asserting optimality, or necessity, of different
multiplication routes for the two initial rows.
Let
\[
 \mathcal T=\{(2,8,1),(3,9,2),(4,10,3),
                (5,8,3),(6,9,4),(12,9,8)\}.
\]
For $N\ge1$ and $r\in\{0,1,2\}$ define $H^*_{N,r}$ to be the product
of primes $p$ with $p^2>50N$ for which, with $t=\lfloor N/p\rfloor$,
at least one $(k,A,B)\in\mathcal T$ satisfies
\[
 \frac{13N+r}{13t+k}<p<\frac{AN}{At+B}.
 \tag{L1}
\]
Every prime occurs once. The inequalities imply $p<8N$, so the product
is finite. Its $t=0$ part is precisely the previous divisor $H_{N,r}$.
If there are no eligible primes, the product is one.

\begin{theorem}
With the same initial rows, trajectory, and factors $C_N,G_N$ as in
the preceding argument, for every integer $N\ge1$ and every
$r\in\{0,1,2\}$,
\[
 \boxed{\frac{C_Nd_{12N}d_{13N+r}}{G_NH^*_{N,r}}
             RW(N)e_r\in\mathbb Z^2.}
\]
\end{theorem}

\begin{proof}
The baseline integer normalization with $G_N$ has already been established.
It suffices to prove one additional factor $p$ for each prime in (L1).
Put $x=N/p=t+f$, where $0\le f<1$. The upper inequality in (L1) says
$f>B/A$; the lower says
\[
 f<\frac{k}{13}-\frac{r}{13p}.
\]
Thus $f$ lies in the same one of the six open intervals as in the
finite chamber table for the original shell argument.
The two intervals with an internal split are divided at $5/11$ and
$10/11$, exactly as before.

Write
\[
 a=\lfloor8x\rfloor=8t+a_0,\quad
 M=\lfloor11x\rfloor=11t+M_0,\quad
 b=8N-ap,\quad d=11N-Mp,\quad h=\phi(x)=\phi(f).
\]
The equality for $h$ follows directly from the periodic floor formula.
For $0\le j<p$, the carry indicators $e=\mathbf1_{j>d}$ and
$\delta_u=\mathbf1_{(j+uN)\bmod p<b}$ depend only on $f$ and $j/p$.
On their minimal-carry support they therefore have the same fixed
values as in the earlier table. The upper digits translate by
\[
 \lambda_u=\lfloor j/p+ux\rfloor=ut+\lambda_{u,0}
 \qquad (u=9,10,11).
\]
In particular they are still independent of the supporting residue $j$.
The two exceptional fractional split points cannot occur: if
$f=5/11$ or $10/11$, primality forces $p=11$ and
$N=11t+5$ or $11t+10$, contrary to $p^2>50N$.

Over $\mathbb F_p$ use the same polynomial quotient
\[
 \mathcal Q(X)=
 \frac{\prod_{u=9,10,11}\binom{X+uN}{b}}
 {(X^p-X)^{h-1}\prod_{j=0}^d(X-j)}.
\]
For $h=0$ the displayed formula means multiplication by $X^p-X$.
The carry multiplicities establish polynomial divisibility exactly as
before. Its decisive degree parameter is now
\[
 K=3a-M+h=13t+(3a_0-M_0+h)=13t+k.
\]
Hence (L1) gives
\[
 \deg(X^r\mathcal Q)=13N+r-(K-1)p-1<p-1,
 \qquad \sum_{j\in\mathbb F_p}j^r\mathcal Q(j)=0.
 \tag{L2}
\]

For clarity, the local kernel argument also remains valid with these
larger upper digits. Set
\[
 P(V)=\prod_{u=9,10,11}(a+1)^{\delta_u}
              \binom{\lambda_u+V}{a+\delta_u},\qquad
 U(V)=\frac{P(V)}{\prod_{l=0}^{M-e}(V-l)},\qquad
 \kappa=(-1)^{d-M+h-1}M!d!.
\]
The cutoff implies $p>50x$; the shell intervals imply $x>1/8$.
Thus $p$ is odd, $a+1<p$, and $M<p$, so all factorial denominators
and distinct kernel poles used here are legitimate over $\mathbb F_p$.
The exact separation of factors divisible by $p$ gives, as previously,
\[
 p^{1-h}R_{\rm ker}(j+pV)
       \equiv\kappa\mathcal Q(j)U(V)\pmod p,
\]
where $R_{\rm ker}(s)=F(s)(-1)^{11N}(11N)!/
\prod_{l=0}^{11N}(s-l)$ is the rational kernel, not the initial-row matrix.
Taking residues supplies the leading coefficients $c_l/p^h$.
Subtracting their partial fractions gives
\[
 p^{1-h}S_r(j+pV)
       \equiv\kappa j^r\mathcal Q(j)[U(V)]_+\pmod p,
 \tag{L3}
\]
with $[U]_+$ its polynomial part. This subtraction includes the
$e=1$ cases and the noncongruent contributions in the original
kernel. Residue classes outside the minimal-carry support give zero.
The previously established integer-valued normalization and degree
bound below $p^2$ justify (L3) also at its removable derivative nodes.

Choose the prime-dependent shift
\[
 \eta=p\left\lceil\frac{2N}{p}\right\rceil=Lp.
\]
If $p\le2N$, then $2N\le\eta<4N$; if $p>2N$, then
$\eta=p<8N$. Thus $2N\le\eta\le9N$ in every case, as required by
the exact readout. In the genuinely new periods $t\ge1$, one has
$p<N$ and the stronger bound $\eta<3N$.
All harmonic indices are below $20N<p^2$ (below $14N$ for $t\ge1$),
and the Newton degree plus one is at most $13N+r<p^2$.
Also $M+L\le13x+1<p$, so the high-digit harmonic reciprocals are units.
The integer factors at the relevant shifted Newton nodes have
absolute values below $24N+r<50N<p^2$.

The harmonic leading residue factors as a constant times the sum in
(L2), because the multiples of $p$ up to $j+(l+L)p$ are indexed by
$l+L$, independently of $j\in\{0,\ldots,p-1\}$.
For the polynomial part, (L3) and Lucas's congruence show that every
Newton coefficient with index $vp-1$ has the same factor
$\sum_jj^r\mathcal Q(j)$ and hence vanishes modulo $p$.
These are the only possibly nonintegral weights in the polynomial
functional. Both contributions therefore have valuation at least
$h-1$; the second row has valuation at least $h$.
Finally $v_p(d_{12N}d_{13N+r})=2$ and $v_p(G_N)=h$, so both
normalized entries acquire the asserted extra $p$.
Multiplying these prime divisibilities proves the statement for all $N$,
including the cases where the eligible product is empty.
\end{proof}

The arithmetic improvement in this theorem comes from the same exact
matrix readout and its prime-dependent polynomial analysis. Separate
finite experiments may expose these factors through specific
multiplication orders; their success alone is not the proof above.
